% Kapittel 2 - Teori
\chapter{Teori}
\label{teori} % Bruk \ref{teori}, for å referere til kapittelet andre steder i oppgaven. 
%----------------------------------------------------------------------------------------
\section{Visual Studio Code}

\section{Angular}
\subsection{Typescript}
\subsection{Node.js}

\section{Chart.js}
\section{Ng2-charts}


\section{Git og Github}

\section{Bootstrap}
Bootstrap er det nest mest fulgte prosjektet på GitHub. Det er et åpen kilde verktøy for å utvikle HTML, CSS og JavaScript. Hovedformålet med å bruke bootstrap er å spare tid på utseendet til applikasjonen ved å bruke klasser og funksjoner som allerede eksisterer i bootstrap. Man får tilgang til knapper, skjema, navigasjonslinjer og andre grensesnittkomponenter\cite{BootstrapAbout}.

\subsection{Tables}
I denne oppgaven ble bootstrap brukt til å utforme tabellen for brushfilene. Det ble laget en enkel tabell i HTML og ".table" klassen fra bootstrap ble lagt til. 
\begin{lstlisting}
<table class="table">
    <thead>
      <tr>
        <th>Fornavn</th>
        <th>Etternavn</th>
      </tr>
    </thead>
    <tbody>
      <tr>
        <td>Ola</td>
        <td>Nordmann</td>
      </tr>
      <tr>
        <td>Ola</td>
        <td>Nordmann</td>
      </tr>
    </tbody>
</table>
\end{lstlisting}
\begin{figure}
    \centering
    \includegraphics{}
    \caption{Enkel HTML-tabell med bootstrap}
    \label{fig:bootstrap_table}
\end{figure}
\cite{BootstrapTable}

\subsection{Grid system}
\cite{BootstrapGrid}






